% !TEX root = ../main.tex

\section{Task B: Read Humidity Values from the HTS221 Sensor}

\subsection{a) Calculations:}

1. The HTS221 must be activated before measurements can be read. For this task, the required startup step is to write
\[
	\texttt{CTRL\_REG1} = \texttt{0x81}
\]
to register \texttt{0x20} via I\textsuperscript{2}C.

The register \texttt{CTRL\_REG1} has the bit structure
\[
	\texttt{PD \;\; 0 \;\; 0 \;\; 0 \;\; 0 \;\; BDU \;\; ODR1 \;\; ODR0}
\]

For the value \texttt{0x81}:
\[
	1000\,0001_2
\]

This means:
\begin{itemize}
	\item \texttt{PD = 1}: sensor active,
	\item \texttt{BDU = 0}: output registers are updated continuously,
	\item \texttt{ODR1:ODR0 = 01}: output data rate \(= 1\,\mathrm{Hz}\).
\end{itemize}

The default value of \texttt{CTRL\_REG1} is
\[
	\texttt{0x00}.
\]
Thus, after reset the sensor is powered down and no periodic measurement is active.

For this task, no further startup configuration is strictly required. Optional registers such as \texttt{AV\_CONF}, \texttt{CTRL\_REG2}, and \texttt{CTRL\_REG3} may remain at their default values if no special averaging, one-shot mode, or data-ready pin is used.

\vspace{0.5cm}

2. The relative humidity is calculated by linear interpolation between two calibration points stored in the sensor.

The measured raw humidity output is:
\[
	H_{\mathrm{OUT}}
\]

The calibration values are:
\[
	H_0,\quad H_1,\quad H_{0T0\_OUT},\quad H_{1T0\_OUT}
\]

The conversion formula is
\[
	RH = H_0 + \frac{H_1 - H_0}{H_{1T0\_OUT} - H_{0T0\_OUT}} \cdot \left(H_{\mathrm{OUT}} - H_{0T0\_OUT}\right)
\]

Equivalently, with the slope
\[
	k = \frac{H_1 - H_0}{H_{1T0\_OUT} - H_{0T0\_OUT}}
\]
the formula can be written as
\[
	RH = k \cdot \left(H_{\mathrm{OUT}} - H_{0T0\_OUT}\right) + H_0
\]

The values \(H_0\) and \(H_1\) are stored in the device as calibration constants. In the registers, they are saved with a factor of 2. Therefore:
\[
	H_0 = \frac{\texttt{H0\_rH\_x2}}{2}, \qquad
	H_1 = \frac{\texttt{H1\_rH\_x2}}{2}
\]

The raw output value is read from two registers:
\[
	H_{\mathrm{OUT}} = (\texttt{HUMIDITY\_OUT\_H} \ll 8)\;|\;\texttt{HUMIDITY\_OUT\_L}
\]

The calibration outputs are also 16-bit signed values:
\[
	H_{0T0\_OUT} = (\texttt{H0\_T0\_OUT\_H} \ll 8)\;|\;\texttt{H0\_T0\_OUT\_L}
\]
\[
	H_{1T0\_OUT} = (\texttt{H1\_T0\_OUT\_H} \ll 8)\;|\;\texttt{H1\_T0\_OUT\_L}
\]

Thus:
\begin{itemize}
	\item \(H_{\mathrm{OUT}}\) changes continuously during measurement,
	\item \(H_0\), \(H_1\), \(H_{0T0\_OUT}\), and \(H_{1T0\_OUT}\) are calibration constants read once from the sensor memory.
\end{itemize}

All required values are found in the HTS221 datasheet and application note in the calibration register section.

\begin{table}[H]
	\centering
	\begin{tabularx}{\textwidth}{|>{\raggedright\arraybackslash}X|>{\raggedright\arraybackslash}X|>{\raggedright\arraybackslash}p{2.7cm}|>{\raggedright\arraybackslash}p{2.2cm}|>{\raggedright\arraybackslash}p{3cm}|}
		\hline
		\textbf{Parameter}        & \textbf{Description}                                      & \textbf{Value or source}                    & \textbf{Unit (if any)} & \textbf{Register address}    \\ \hline
		\texttt{CTRL\_REG1}       & Startup configuration register                            & set to \texttt{0x81}                        & --                     & \texttt{0x20}                \\ \hline
		\texttt{H0\_rH\_x2}       & First humidity calibration point, stored multiplied by 2  & sensor calibration constant                 & \% rH $\times 2$       & \texttt{0x30}                \\ \hline
		\texttt{H1\_rH\_x2}       & Second humidity calibration point, stored multiplied by 2 & sensor calibration constant                 & \% rH $\times 2$       & \texttt{0x31}                \\ \hline
		\(H_0\)                   & First humidity calibration value                          & \(\texttt{H0\_rH\_x2}/2\)                   & \% rH                  & derived from \texttt{0x30}   \\ \hline
		\(H_1\)                   & Second humidity calibration value                         & \(\texttt{H1\_rH\_x2}/2\)                   & \% rH                  & derived from \texttt{0x31}   \\ \hline
		\texttt{H0\_T0\_OUT\_L}   & Low byte of first calibration output                      & sensor calibration constant                 & --                     & \texttt{0x36}                \\ \hline
		\texttt{H0\_T0\_OUT\_H}   & High byte of first calibration output                     & sensor calibration constant                 & --                     & \texttt{0x37}                \\ \hline
		\(H_{0T0\_OUT}\)          & First raw calibration output                              & \((\texttt{0x37} \ll 8)\;|\;\texttt{0x36}\) & LSB                    & \texttt{0x36}, \texttt{0x37} \\ \hline
		\texttt{H1\_T0\_OUT\_L}   & Low byte of second calibration output                     & sensor calibration constant                 & --                     & \texttt{0x3A}                \\ \hline
		\texttt{H1\_T0\_OUT\_H}   & High byte of second calibration output                    & sensor calibration constant                 & --                     & \texttt{0x3B}                \\ \hline
		\(H_{1T0\_OUT}\)          & Second raw calibration output                             & \((\texttt{0x3B} \ll 8)\;|\;\texttt{0x3A}\) & LSB                    & \texttt{0x3A}, \texttt{0x3B} \\ \hline
		\texttt{HUMIDITY\_OUT\_L} & Low byte of measured humidity output                      & current sensor value                        & --                     & \texttt{0x28}                \\ \hline
		\texttt{HUMIDITY\_OUT\_H} & High byte of measured humidity output                     & current sensor value                        & --                     & \texttt{0x29}                \\ \hline
		\(H_{\mathrm{OUT}}\)      & Current raw humidity output                               & \((\texttt{0x29} \ll 8)\;|\;\texttt{0x28}\) & LSB                    & \texttt{0x28}, \texttt{0x29} \\ \hline
		\(RH\)                    & Calculated relative humidity                              & computed from interpolation formula         & \% rH                  & --                           \\ \hline
	\end{tabularx}
\end{table}

\subsection{b) Implementation:}

\textit{Describe your solution. Add all relevant code snippets into a listing or as screenshots and describe their purpose. Remember to use generic methods with well-defined parameters!!}

\vspace{3cm}

\subsection{c) Results:}

\textit{Add a screenshot of the UART output.}

\vspace{4cm}

\subsection{d) Discussion:}

\textit{Describe your experiences (e.g., design decisions, problems, lesson learned). Which part of the code will be reusable?}

\vspace{2cm}
